Mathematical discovery has historically been a human activity in which abstraction, search, and validation are interleaved in the mind of the mathematician. Modern tools have externalised each of these stages in different communities. Mechanised theorem provers such as Lean and Isabelle externalise validation. Large language models trained on mathematical corpora externalise informal abstraction and proof sketching. Quantum algorithms such as Grover amplification externalise unstructured search with a quadratic speedup over classical brute force. Causal inference machinery based on do-interventions externalises counterfactual reasoning. Mechanistic interpretability through circuit probes externalises the inspection of the internal computations of a model.

No existing system brings all of these externalisations together in a single closed loop. The Large Reasoning and Action Model (LRAM) introduced in this whitepaper is an attempt to do exactly that. The motivating question is whether such an integration can be more than the sum of its parts for the hardest mathematical problems of our era, the Clay Millennium Prize Problems, which include the Riemann Hypothesis, the P versus NP question, the existence and smoothness of Navier-Stokes solutions, the Yang-Mills mass gap, the Hodge Conjecture, the Birch and Swinnerton-Dyer Conjecture, and the resolved but methodologically instructive Poincare Conjecture.

The purpose of this document is not to claim a solution to these problems. The purpose is to state an honest architecture, to derive its formal properties, to bound its limits, to implement it in open-source tooling, and to publish a reproducible proof-of-value pipeline that future researchers can extend. We apply a disciplined set of epistemic tools throughout: first principles decomposition, algorithmic thinking, strong mathematical foundations, higher-dimensional lifting, the Paul-Elder critical-thinking framework, clarity, inversion, and Reflexion-enhanced self-critique.

The remainder of the whitepaper is organised as follows. Sections 3 and 4 establish the epistemic and first-principles foundation. Section 5 is the formal heart of the architecture. Sections 6 and 7 extend the formalism to causal reasoning, world-model distillation, and teacher-student learning. Sections 8, 9 and 10 develop inversion analysis, higher-dimensional lifting, and the Reflexion loop. Sections 11 and 12 describe the reproducible workshop and interpret its results. Sections 13 and 14 bound the programme by stating its limits and mapping a staged roadmap. Section 15 concludes.
